\documentclass{beamer}
\usetheme{default}
\usepackage{amsmath,amssymb}

\title{Routing System Optimization}
\author{}
\date{}

\begin{document}

\frame{\titlepage}

\begin{frame}{Problem Setup}

We consider a routing system with $K$ LLMs.

\textbf{Prompt distribution}

\[
X \sim \mathcal{D}
\]

\textbf{Routing policy}

\[
M = \pi(X)
\]

Even if $\pi$ is deterministic, $M$ is random because $X$ is random.

\textbf{Routing fractions}

\[
w_i = \Pr(M=i)
\]

with

\[
w \in \Delta_K
\]

where $\Delta_K$ is the probability simplex.

\end{frame}

\begin{frame}{Score Model}

For each prompt $X$, each model $i$ produces a score

\[
S(X) = (s_1(X), \dots, s_K(X))
\]

Under routing policy $\pi$:

\[
S_\pi(X) = s_{\pi(X)}(X)
\]

The expected score becomes

\[
\mathbb{E}[S_\pi]
=
\sum_{i=1}^K
w_i
\mathbb{E}[s_i(X) \mid M=i]
\]

Define

\[
S(w)
\]

as the \textbf{maximum achievable score} under routing fractions $w$.

\end{frame}

\begin{frame}{Latency Model}

Let the global request rate be

\[
\lambda
\]

If routing fraction is $w_i$, model $i$ receives load

\[
\lambda_i = \lambda w_i
\]

For setup $\Theta$, the expected latency of model $i$ is

\[
\ell_i(\theta_i, \lambda_i)
\]

The overall expected latency becomes

\[
L(\Theta, w)
=
\sum_{i=1}^{K}
w_i
\,
\ell_i(\theta_i, \lambda w_i)
\]

We estimate $\ell_i$ using a learned latency predictor.

\end{frame}

\begin{frame}{Offline Optimization Problem}

Our goal is to select

\[
(\Theta, w)
\]

to maximize score while satisfying latency constraints.

\[
\max_{\Theta,\, w \in \Delta_K}
S(w)
\]

subject to

\[
L(\Theta, w) \le \tau
\]

where $\tau$ is the latency SLO.

\end{frame}

\begin{frame}{Lagrangian Formulation}

We introduce a Lagrange multiplier $\beta$.

\[
\mathcal{L}(w, \beta)
=
S(w)
-
\beta \left(L(\Theta, w) - \tau\right)
\]

The constrained problem becomes

\[
\min_{\beta \ge 0}
\max_{w \in \Delta_K}
\mathcal{L}(w,\beta)
\]

This converts the constrained problem into an unconstrained saddle-point problem.

\end{frame}

\begin{frame}{Gradient of the Score Function}

The optimal score under routing fractions is

\[
S(w)
=
\min_{\alpha}
\left[
\frac{1}{N}
\sum_{j=1}^N
\max_i (s_{ji} - \alpha_i)
+
\sum_{i=1}^K
w_i \alpha_i
\right]
\]

Let $\alpha^*(w)$ be the optimal dual prices.

Using the envelope theorem,

\[
\nabla S(w) = \alpha^*(w)
\]

Thus the dual prices directly give the gradient of the optimal score.

\end{frame}

\begin{frame}{Optimization Algorithm}

We perform projected gradient ascent on $w$.

Update rule:

\[
w \leftarrow
\Pi_{\Delta_K}
\Big(
w + \eta
(
\alpha^*(w)
-
\beta
\nabla L(\Theta, w)
)
\Big)
\]

where

\begin{itemize}
\item $\alpha^*(w)$ is the dual price vector
\item $\nabla L(\Theta,w)$ is obtained from the latency model
\item $\Pi_{\Delta_K}$ projects onto the simplex
\end{itemize}

\end{frame}

\begin{frame}{Key Insight}

The routing optimizer balances:

\begin{itemize}
\item score improvement from routing more prompts to a model
\item latency penalty due to increased load
\end{itemize}

The dual prices $\alpha^*(w)$ represent the marginal value of assigning more prompts to each model.

\end{frame}

\end{document}
